\documentclass[a4paper,12pt]{article}
\usepackage[T2A]{fontenc}
\usepackage[utf8]{inputenc}
\usepackage[russian]{babel}
\usepackage{amsmath,amssymb,mathtools}
\usepackage{helvet}
\renewcommand{\sfdefault}{cmss}
\renewcommand{\familydefault}{\sfdefault}
\usepackage{icomma}
\usepackage[margin=23mm,headheight=23pt,footskip=16mm]{geometry}
\usepackage{microtype}
\usepackage{tikz}
\usetikzlibrary{calc,arrows.meta,positioning,shapes.geometric}
\usepackage{enumitem}
\usepackage{tabularx,booktabs}
\usepackage{fancyhdr}
\usepackage{setspace}
\usepackage{xcolor}
\usepackage{hyperref}

\definecolor{accent}{HTML}{087B7B}
\definecolor{darkaccent}{HTML}{075455}
\definecolor{textgray}{HTML}{343A3D}
\definecolor{softgray}{HTML}{707D80}
\definecolor{fine}{HTML}{D5DFDF}
\hypersetup{colorlinks=true,linkcolor=darkaccent,
 pdfauthor={Лёвин Артём Александрович},
 pdftitle={Чек-лист. Метод интервалов, рациональные уравнения, метод Гаусса}}
\setstretch{1.11}
\setlength{\parindent}{0pt}
\setlength{\parskip}{0.44em}
\setlist[enumerate]{leftmargin=*,itemsep=0.48em,topsep=0.45em}
\pagestyle{fancy}
\fancyhf{}
\fancyhead[L]{\footnotesize\color{softgray} Математика \,\textbullet\, 10 класс}
\fancyhead[R]{\footnotesize\color{softgray} 20.10.26}
\fancyfoot[C]{\footnotesize\color{softgray}\thepage}
\renewcommand{\headrulewidth}{0.35pt}
\renewcommand{\headrule}{\hbox to\headwidth{\color{fine}\leaders\hrule height \headrulewidth\hfill}}
\newcommand{\subnote}[1]{{\color{softgray}\small #1}}
\newcommand{\strong}[1]{{\color{darkaccent}\textbf{#1}}}
\newcommand{\ans}[1]{\textbf{#1}}
\tikzset{
  flowstart/.style={ellipse,draw=darkaccent,line width=.9pt,align=center,
    minimum width=31mm,minimum height=11mm,fill=white,font=\small},
  flowact/.style={rounded corners=2pt,rectangle,draw=accent,line width=.9pt,
    align=center,text width=100mm,minimum height=13mm,inner sep=5pt,font=\small},
  flowhalf/.style={rounded corners=2pt,rectangle,draw=accent,line width=.9pt,
    align=center,text width=45mm,minimum height=14mm,inner sep=4pt,font=\small},
  flowcond/.style={diamond,aspect=2.7,draw=darkaccent,line width=.9pt,
    align=center,inner sep=3pt,text width=39mm,font=\small},
  arr/.style={-{Latex[length=3mm,width=2mm]},thick,draw=darkaccent}
}

\begin{document}
\thispagestyle{empty}
\vspace*{32mm}
\begin{center}
{\color{accent}\large МАТЕМАТИКА \,\textbullet\, 10 КЛАСС}\\[14mm]
{\color{softgray}\large Персональный учебный материал}\\[8mm]
{\Huge\bfseries Чек-лист}\\[9mm]
{\LARGE\bfseries\color{darkaccent} Метод интервалов\\[2mm]
Рациональные уравнения\\[2mm] Метод Гаусса}\\[17mm]
{\large Повторение ключевых приёмов}\\[18mm]
{\color{textgray}\large 20 октября 2026 года}
\end{center}
\vfill
\begin{center}
{\color{darkaccent}\large Лёвин Артём Александрович}\\[2mm]
{\large эксклюзивно для Анны Трапезниковой}
\end{center}
\vspace{14mm}
\begin{tikzpicture}[remember picture,overlay]
\draw[accent!60,line width=1pt]
([xshift=23mm,yshift=19mm]current page.south west)
.. controls ([xshift=66mm,yshift=33mm]current page.south west)
and ([xshift=-64mm,yshift=5mm]current page.south east)
.. ([xshift=-23mm,yshift=20mm]current page.south east);
\end{tikzpicture}
\clearpage

\section*{\color{darkaccent}01.\quad Метод интервалов и кратность корня}
\strong{Главная мысль.} На числовой прямой знак многочлена может
измениться только при переходе через его корень. При переходе через корень
\textbf{нечётной кратности} знак меняется, а \textbf{чётной кратности}
сохраняется.

\begin{center}
\begin{tabularx}{\linewidth}{@{}p{46mm}X@{}}
\toprule
Вид множителя & Что происходит со знаком при переходе через корень\\
\midrule
$(x-a)^{2k+1}$ & Знак меняется.\\
$(x-a)^{2k}$ & Знак сохраняется.\\
\bottomrule
\end{tabularx}
\end{center}

\textbf{Пример из занятия:}
\[
 x^2(5x-4)(x+7)<0.
\]
Нули: $x=-7$, $x=0$ (кратность $2$), $x=\dfrac45$.
Справа от $\frac45$ знак положительный: при $x=1$ получаем $8>0$.
Переходя справа налево, меняем знак на $\frac45$, сохраняем при $0$
и меняем на $-7$.

\begin{center}
\begin{tikzpicture}[x=1cm,y=1cm]
\draw[-{Latex[length=2.5mm]},thick] (0,0)--(12.6,0);
\foreach \x/\lab in {2/{-7},6/{0},9/{\frac45}} {
\draw[fill=white,line width=1.15pt,draw=darkaccent] (\x,0) circle (2.7pt);
\node[below=5pt,font=\small] at (\x,0) {$\lab$};
}
\node[above=6pt] at (1,0) {$+$};
\node[above=6pt] at (4,0) {$-$};
\node[above=6pt] at (7.5,0) {$-$};
\node[above=6pt] at (11,0) {$+$};
\draw[accent,line width=2.4pt] (2.10,0)--(5.90,0);
\draw[accent,line width=2.4pt] (6.10,0)--(8.90,0);
\end{tikzpicture}
\end{center}
\[
\boxed{x\in(-7;0)\cup\left(0;\frac45\right).}
\]
\strong{Почему ноль выколот?} Хотя знак при переходе через $x=0$
сохраняется, сама левая часть равна нулю. В строгом неравенстве
нулевое значение запрещено.

\textbf{Строгость и объединение интервалов.}
Для $F(x)=(x-4)^2(x+5)^3$ получаем:
\[
\begin{array}{@{}ll@{}}
F(x)>0: & (-5;4)\cup(4;+\infty),\\[1mm]
F(x)\ge0: & [-5;+\infty),\\[1mm]
F(x)<0: & (-\infty;-5),\\[1mm]
F(x)\le0: & (-\infty;-5]\cup\{4\}.
\end{array}
\]
Точку $4$ можно «склеить» с соседними промежутками \emph{только тогда},
когда она входит в решение.

\clearpage
\section*{\color{darkaccent}Алгоритм 1.\quad Как решить неравенство методом интервалов}
\vspace{4mm}
\begin{center}
\begin{tikzpicture}[node distance=10mm]
\node[flowstart] (start) {Начало};
\node[flowact,below=of start] (roots) {Приравнять левую часть к нулю.\\ Найти корни и их кратности};
\node[flowact,below=of roots] (dots) {Отметить корни на оси, упорядочив их.\\ Отдельно учесть ОДЗ, если есть знаменатель};
\node[flowact,below=of dots] (test) {Определить знак на одном промежутке\\ (например, подстановкой удобного числа)};
\node[flowcond,below=12mm of test] (cond) {Кратность очередного корня чётная?};
\node[flowhalf] (even) at ($(cond.south)+(-4.0cm,-2.8cm)$) {При переходе\\ сохранить знак};
\node[flowhalf] (odd) at ($(cond.south)+(4.0cm,-2.8cm)$) {При переходе\\ поменять знак};
\node[flowact] (select) at ($(cond.south)+(0cm,-5.5cm)$)
{Проставить знаки на всех промежутках.\\ Выбрать промежутки нужного знака};
\node[flowact,below=of select] (ends) {Проверить строгость и ОДЗ;\\ включить или исключить граничные точки};
\node[flowstart,below=of ends] (finish) {Записать ответ};
\draw[arr] (start)--(roots);
\draw[arr] (roots)--(dots);
\draw[arr] (dots)--(test);
\draw[arr] (test)--(cond);
\draw[arr] (cond.west)-| node[pos=.24,above,font=\small]{да} (even.north);
\draw[arr] (cond.east)-| node[pos=.24,above,font=\small]{нет} (odd.north);
\draw[arr] (even.south) |- (select.west);
\draw[arr] (odd.south) |- (select.east);
\draw[arr] (select)--(ends);
\draw[arr] (ends)--(finish);
\end{tikzpicture}
\end{center}
\clearpage

\section*{\color{darkaccent}02.\quad Рациональные уравнения: замена переменной}
\strong{Когда искать замену?} Если в разных частях уравнения
повторяется одно и то же выражение или встречаются взаимно обратные дроби.
Обозначайте через $t$ такой фрагмент, при котором новое уравнение
становится проще.

\textbf{Пример 1. Выделяем весь знаменатель:}
\[
\frac{12}{x^2-2x-3}-(x^2-2x)=-1.
\]
\strong{ОДЗ:} $x^2-2x-3\ne0$, то есть $x\ne-1$, $x\ne3$.
Вводим $t=x^2-2x-3$, причём $t\ne0$. Тогда $x^2-2x=t+3$, поэтому
\[
\frac{12}{t}-(t+3)=-1
\quad\Longleftrightarrow\quad
\frac{12}{t}=t+2
\quad\Longleftrightarrow\quad
t^2+2t-12=0.
\]
\[
t=-1\pm\sqrt{13}.
\]
Возвращаемся к исходной переменной:
\[
x^2-2x-3=t\quad\Longleftrightarrow\quad (x-1)^2=t+4.
\]
Для $t=-1-\sqrt{13}$ имеем $t+4=3-\sqrt{13}<0$,
поэтому действительных $x$ нет. Для $t=-1+\sqrt{13}$ получаем
\[
\boxed{x=1\pm\sqrt{3+\sqrt{13}}.}
\]
Оба корня удовлетворяют исходному ОДЗ.

\textbf{Проверочный список}
\begin{enumerate}[label=$\square$,leftmargin=10mm,itemsep=.29em]
\item Записал ограничения знаменателей \emph{до} преобразований.
\item Выбрал замену, упрощающую дроби.
\item Все повторяющиеся выражения выразил через $t$.
\item Решил уравнение относительно $t$ и отбросил недопустимые $t$.
\item Выполнил обратную замену и проверил допустимость $x$.
\end{enumerate}

\clearpage
\section*{\color{darkaccent}03.\quad Взаимно обратные выражения}
\textbf{Пример 2.} Рассмотрим уравнение
\[
\frac{4(x+5)}{x^2+5x+14}+\frac{x^2+5x+14}{x+5}=5.
\]
\strong{ОДЗ:} $x\ne-5$. Квадратный трёхчлен $x^2+5x+14$
не обращается в ноль при действительных $x$, поскольку его
дискриминант $25-56=-31<0$.

Пусть $t=\dfrac{x+5}{x^2+5x+14}$, $t\ne0$. Вторая дробь равна $1/t$:
\[
4t+\frac1t=5\quad\Longleftrightarrow\quad
4t^2-5t+1=0
\quad\Longleftrightarrow\quad t=1\ \text{или}\ t=\frac14.
\]
При $t=1$ получаем $x^2+4x+9=0$ --- действительных корней нет.
При $t=\frac14$ имеем $x^2+x-6=0$, откуда
\[
\boxed{x=-3\ \text{или}\ x=2.}
\]
\subnote{Обратите внимание: при замене дроби вида $A/B$ обратная дробь $B/A$
становится $1/t$. Важно, чтобы $A\ne0$ и $B\ne0$.}

\medskip
\section*{\color{darkaccent}04.\quad Разность квадратов в уравнениях}
Если в уравнении можно увидеть два квадрата,
используем формулу
\[
A^2-B^2=(A-B)(A+B).
\]
Таким образом, $A^2=B^2$ равносильно совокупности $A=B$ или $A=-B$
при соблюдении всех исходных ограничений.

\textbf{Короткий пример:}
\[
\left(\frac{x+1}{x-2}\right)^2=4,
\qquad x\ne2.
\]
Замечаем, что $4=2^2$, переносим вправо и раскладываем разность
квадратов:
\[
\left(\frac{x+1}{x-2}-2\right)
\left(\frac{x+1}{x-2}+2\right)=0.
\]
Отсюда либо $x+1=2(x-2)$, либо $x+1=-2(x-2)$.
Получаем $x=5$ или $x=1$. Оба значения допустимы.

\strong{Типичная ошибка.} Чтобы применять разность квадратов,
\emph{оба} слагаемых должны иметь вид квадратов. Например,
$2=(\sqrt2)^2$ и $4=2^2$.

\clearpage
\section*{\color{darkaccent}05.\quad Метод Гаусса: повторение}
Записываем систему расширенной матрицей; строки обозначаем $R_1,R_2,R_3$.
Последовательно обнуляем элементы \emph{под} ведущими коэффициентами,
затем выполняем обратный ход.

\textbf{Пример с отсутствующим $z$ в третьем уравнении:}
\[
\begin{cases}
x+y+z=6,\\
2x-y+z=3,\\
3x-y=1.
\end{cases}
\quad\Longleftrightarrow\quad
\left(\begin{array}{rrr|r}
1&1&1&6\\2&-1&1&3\\3&-1&0&1
\end{array}\right).
\]
Здесь отсутствие $z$ означает коэффициент $\mathbf{0}$.
\[
\begin{aligned}
R_2&\leftarrow R_2-2R_1,\\
R_3&\leftarrow R_3-3R_1
\end{aligned}
\quad\Longrightarrow\quad
\left(\begin{array}{rrr|r}
1&1&1&6\\0&-3&-1&-9\\0&-4&-3&-17
\end{array}\right).
\]
\[
R_3\leftarrow 3R_3-4R_2
\quad\Longrightarrow\quad
\left(\begin{array}{rrr|r}
1&1&1&6\\0&-3&-1&-9\\0&0&-5&-15
\end{array}\right).
\]
Обратный ход: $z=3$, затем $-3y-z=-9\Rightarrow y=2$,
затем $x+y+z=6\Rightarrow x=1$.
Итак, $(x,y,z)=(1,2,3)$.

\strong{Как быстро увидеть отсутствие решений?}
Если преобразования привели к строке
\[
(0\quad 0\quad 0\mid c),\qquad c\ne0,
\]
она означает противоречие $0=c$. \textbf{Решений нет.}
Напротив, строка $(0\ 0\ 0\mid0)$ означает тождество;
в этом случае проверяют число свободных переменных.

\clearpage
\section*{\color{darkaccent}Алгоритм 2.\quad Прямой и обратный ход Гаусса}
\vspace{5mm}
\begin{center}
\begin{tikzpicture}[node distance=9mm]
\node[flowstart] (a) {Начало};
\node[flowact,below=of a] (b) {Записать расширенную матрицу\\ (отсутствующий коэффициент заменять нулём)};
\node[flowact,below=of b] (c) {Выбрать ненулевой ведущий элемент.\\ При необходимости переставить строки};
\node[flowact,below=of c] (d) {Обнулить элементы ниже ведущего.\\ Преобразовывать строки целиком};
\node[flowact,below=of d] (e) {Повторить для следующих столбцов.\\ Получить ступенчатую матрицу};
\node[flowcond,below=12mm of e] (f) {Есть строка $0=c$, где $c\ne0$?};
\node[flowhalf,right=19mm of f] (g) {Система несовместна.\\ Решений нет};
\node[flowact,below=19mm of f] (h) {Если противоречия нет --- обратный ход.\\ Найти ведущие неизвестные};
\node[flowstart,below=of h] (i) {Проверить решение};
\draw[arr] (a)--(b);
\draw[arr] (b)--(c);
\draw[arr] (c)--(d);
\draw[arr] (d)--(e);
\draw[arr] (e)--(f);
\draw[arr] (f.east) -- node[above,font=\small]{да} (g.west);
\draw[arr] (f.south) -- node[right,font=\small]{нет} (h.north);
\draw[arr] (h)--(i);
\end{tikzpicture}
\end{center}
\clearpage

\section*{\color{darkaccent}06.\quad Самостоятельная тренировка}
\textbf{Путь А \,\textbullet\, 10 заданий}

\subnote{В № 1--4 решите неравенства методом интервалов;
в № 5--7 используйте замену переменной; в № 8 --- разность квадратов;
в № 9--10 примените метод Гаусса. В уравнениях сначала указывайте ОДЗ.}

\subsection*{\color{darkaccent}Часть I. Метод интервалов}
\begin{enumerate}[label=\textbf{\arabic*.},itemsep=.75em]
\item $\displaystyle (x-2)^2(x+3)>0.$
\item $\displaystyle (x+1)^2(x-4)\le0.$
\item $\displaystyle x^2(2x-3)(x+2)<0.$
\item $\displaystyle (x-1)^2(x+2)^3\ge0.$
\end{enumerate}

\subsection*{\color{darkaccent}Часть II. Рациональные уравнения}
\begin{enumerate}[label=\textbf{\arabic*.},start=5,itemsep=1em]
\item $\displaystyle \frac6{x^2-3x+2}=(x^2-3x+2)+1.$
\item $\displaystyle \frac{2(x+1)}{x^2+x+1}+\frac{x^2+x+1}{x+1}=3.$
\item $\displaystyle \frac{x^2+2x-4}{x^2+2x-1}
+\frac{x^2+2x-1}{x^2+2x+2}=2.$
\item $\displaystyle \left(\frac{x+2}{x-1}\right)^2=9.$
\end{enumerate}

\subsection*{\color{darkaccent}Часть III. Метод Гаусса}
\begin{enumerate}[label=\textbf{\arabic*.},start=9,itemsep=1em]
\item $\displaystyle \begin{cases}
x+y+z=6,\\ 2x-y+z=3,\\ 3x-y=1.
\end{cases}$
\item $\displaystyle \begin{cases}
x+y+z=4,\\ 2x-y+z=1,\\ 3x+2z=7.
\end{cases}$
\end{enumerate}

\clearpage
\section*{\color{darkaccent}07.\quad Ответы для самопроверки}
\subnote{Сначала завершите решение самостоятельно; затем сверьте итог.}\par\vspace{5mm}
\renewcommand{\arraystretch}{1.6}
\begin{tabularx}{\linewidth}{@{}p{14mm}X@{}}
\toprule
№ & Ответ\\
\midrule
1 & $x\in(-3;2)\cup(2;+\infty)$.\\
2 & $x\in(-\infty;4]$.\\
3 & $x\in(-2;0)\cup\left(0;\dfrac32\right)$.\\
4 & $x\in[-2;+\infty)$.\\
5 & $x=0$ или $x=3$.\\
6 & $x=0$ или $\displaystyle x=\frac{1-\sqrt5}{2}$ или
$\displaystyle x=\frac{1+\sqrt5}{2}$.\\
7 & $\displaystyle x=-1\pm\frac{\sqrt2}{2}$.\\
8 & $\displaystyle x=\frac14$ или $x=\frac52$.\\
9 & $(x,y,z)=(1,2,3)$.\\
10 & Решений нет.\\
\bottomrule
\end{tabularx}

\vspace{8mm}
\textbf{Контрольные ориентиры}
\begin{itemize}[leftmargin=7mm,itemsep=.55em]
\item В № 1 точка $x=2$ выколота: неравенство строгое.
\item В № 4 точка $x=1$ уже включена в непрерывный промежуток
$[-2;+\infty)$, отдельно выделять её не требуется.
\item В № 5 при замене $t=x^2-3x+2$ значения $t=0$ запрещены;
кандидат $t=-3$ не приводит к действительным $x$.
\item В № 7 удобна замена $t=x^2+2x-1$; учитывайте $t\ne0$ и $t\ne-3$.
\item В № 10 сумма первых двух уравнений даёт $3x+2z=5$,
хотя третье требует $3x+2z=7$.
\end{itemize}
\end{document}
